% GEP Intern -- Technology | How R1 + R2 + HR actually go
% For Ojas Srivastava --- speak only what is on the submitted resume
\documentclass[a4paper,9.5pt]{extarticle}
\usepackage[margin=0.85cm]{geometry}
\usepackage{enumitem}
\usepackage{titlesec}
\usepackage{xcolor}
\usepackage{tabularx}
\usepackage{array}
\usepackage{amssymb}
\usepackage[T1]{fontenc}

\pagestyle{plain}
\definecolor{ink}{HTML}{111827}
\definecolor{muted}{HTML}{4B5563}
\definecolor{accent}{HTML}{0F766E}
\definecolor{hot}{HTML}{B91C1C}
\definecolor{ok}{HTML}{166534}
\color{ink}
\setlength{\parskip}{1.5pt}
\setlist[itemize]{leftmargin=1.1em, topsep=1pt, itemsep=1.4pt, parsep=0pt}
\titleformat{\section}{\bfseries\large\color{accent}}{}{0em}{}[\vspace{0.5pt}{\color{accent}\hrule height 0.55pt}]
\titlespacing{\section}{0pt}{8pt}{3pt}
\titleformat{\subsection}{\bfseries\normalsize}{}{0em}{}
\titlespacing{\subsection}{0pt}{5pt}{1.5pt}

\newcommand{\Q}[1]{\par\vspace{1.5pt}\noindent\textcolor{hot}{\textbf{Q:} #1}\par}
\newcommand{\Typ}[1]{\noindent\textcolor{muted}{\textbf{What people usually say:} #1}\par}
\newcommand{\Say}[1]{\noindent\textcolor{ok}{\textbf{You say:} #1}\par}

\begin{document}
\begin{center}
  {\LARGE\textbf{GEP Intern -- Technology: Interview Playbook}}\\[2pt]
  {\small\color{hot}\textbf{How R1 (technical) and R2 (technical / functional + HR) actually go --- questions + answers}}\\[3pt]
  {\small Ojas Srivastava \textbar{} B.Tech AI, SVNIT Surat \textbar{} Submitted resume freeze \textbar{} Hyderabad / Mumbai}\\[1pt]
  {\footnotesize\color{muted} Built from GEP campus reports (GFG, LinkedIn intern offers, AmbitionBox, GEP Campus Connect) + your resume. Not a generic DSA dump.}
\end{center}

\section{0. How the loop really works}
GEP campus tech hiring is consistent across intern and SDE reports:
\textbf{OA} $\to$ \textbf{Technical 1} $\to$ \textbf{Technical 2} $\to$ \textbf{HR} (sometimes folded into R2).
Campus drives often finish in 1--2 days. Interviewers \textbf{do not use a fixed question paper}. They take your intro + resume and drill whatever you said. Overstatement is the main reject reason.

\noindent\begin{tabularx}{\textwidth}{@{}>{\bfseries}p{2.6cm} X@{}}
R1 ($\sim$45--60 min) & Intro $\to$ OOP / CS fundamentals $\to$ 1 Easy--Medium DSA (arrays, hashing, graphs/trees) $\to$ one project deep-dive $\to$ sometimes a puzzle. \\
R2 ($\sim$35--50 min) & Deeper on \textbf{the same projects} + more CS / a second coding question (Fibonacci in a range, reverse array) + ``design this using OOP'' + AI/LLM if you mentioned it. Functional: can you explain LogiFlow like a product. \\
HR ($\sim$10--20 min) & Why GEP, family, location (Hyd \emph{and} Mumbai --- stay open), strengths/weakness, any questions. Not a formality if you sound uninterested or rigid on location. \\
\end{tabularx}

\textbf{Pattern from people who got intern/SDE offers:}
foundations $>$ hard CP; project ownership $>$ buzzwords; think-aloud $>$ silent coding.
A GEP intern (Shivam Kumar, campus SDE 2025): one Easy--Medium DSA per technical round (arrays/graphs) + long project discussion; questions came from \emph{his} intro and code, not a bank.
A GEP intern from SVNIT (Sarthak Chauhan): R1 OOP pillars + project auth; R2 favourite DS (trees) then BFS-style problem + puzzle; HR hobbies + \emph{can you relocate}.

\textbf{Your drive (from PPT):} two interview rounds after OA. Treat R1 as fundamentals + LogiFlow. Treat R2 as Community Hero / Career Automation / Nexus + HR. IFFCO is on the resume --- they can still ask even if you skip it in the intro. \textbf{AirHelp is not on the submitted resume} --- old cheat sheets still mention it; do not volunteer it.

\section{1. Locked intro (speak this, then stop)}
\textit{Good afternoon sir, ma'am. I am Ojas Srivastava. I am a pre-final year B.Tech student in Artificial Intelligence at SVNIT Surat, with a CGPA of 9.20.}

\textit{My coursework includes Object-Oriented Programming, Database Management Systems, Operating Systems, Computer Networks, and Design and Analysis of Algorithms.}

\textit{I was part of the team that secured Global rank 106 in Google Solution Challenge 2026. We built LogiFlow, a supply-chain routing system. I also built Community Hero for Google Vibe2Ship 2026, Global Top 20 --- a civic issue reporting PWA.}

\textit{My coding workflow is built around Cursor, with MCP and plugins. I still write the spec and review every merge.}

\textit{Currently I am Chairperson of Nexus, SVNIT --- the departmental cell of CSE and AI. We manage and mentor 500-plus students.}

\textit{I treat that the same way I treat a project: I own the outcome. GEP works on AI-powered procurement and supply chain. That is the kind of work I have already started on --- taking operational data to a decision, and owning it till it ships.}

\vspace{2pt}
\noindent\textbf{Do not say in R1:} Intern--Technology title, Hyderabad/Mumbai, IFFCO (unless they open the resume), MCP product names, LangGraph, Azure, Angular, ``main lead of the whole team.''\\
\textbf{Do say if asked:} technical lead on railway / comparator / hybrid pipelines; ranking on time, cost, delay risk --- not ``type of goods'' unless you can defend it from the code.

\section{2. Minute-by-minute: Technical R1}
\begin{enumerate}[leftmargin=1.3em, itemsep=2pt]
  \item \textbf{0--2 min --- camera, resume, ``tell me about yourself.''} Speak the locked intro. Stop. Smile. Do not add IFFCO or location.
  \item \textbf{2--12 min --- they pick a word from the intro.} Most likely LogiFlow or OOP from coursework. Have a 60-second LogiFlow story ready (section 4).
  \item \textbf{12--22 min --- CS fundamentals.} OOP 10 minutes is extremely common. Then OS or DBMS if time.
  \item \textbf{22--40 min --- one DSA.} Clarify $\to$ brute $\to$ better $\to$ code $\to$ dry-run $\to$ complexity. Never silent.
  \item \textbf{40--50 min --- puzzle or ``any questions?''} Ask one technical question (section 9). Do not ask stipend here.
\end{enumerate}

\section{3. Minute-by-minute: Technical / functional R2 + HR}
R2 is not a second copy of R1. Reports: they go \textbf{deeper on one project} (ML/AI if you have it --- you do: Gemini on Community Hero; delay ranking on LogiFlow), then a \textbf{simpler coding} question, then culture.

\begin{enumerate}[leftmargin=1.3em, itemsep=2pt]
  \item \textbf{0--3 min --- intro again?} 30-second version: name, LogiFlow Top 106, Community Hero Top 20, Nexus chairperson, GEP close. Do not repeat full coursework.
  \item \textbf{3--20 min --- functional:} ``Explain LogiFlow to a product manager.'' Time, cost, delay. Who uses it. What breaks.
  \item \textbf{20--35 min --- AI they heard in PPT:} What is an LLM vs a rules engine? What is a guardrail? Do \emph{not} invent LangGraph.
  \item \textbf{35--45 min --- second code or OOP design} (bank / cart / Zoomcar-style).
  \item \textbf{HR block:} why GEP, location \textbf{both}, family, weakness, joining date. Stay open to Hyderabad and Mumbai.
\end{enumerate}

\textbf{Intern-length warning:} one on-campus intern R2 was $\sim$10 min (OOP purpose, DS vs algorithms, a lateral-thinking web Q, ``where in 5 years''). AI-intern campuses have used names \textbf{Preliminary} then \textbf{Functional} instead of Tech 1 / Tech 2 --- same content. Do not panic if R2 is short; still finish cleanly.

Official: process aims to close in 1--2 days. Fail a campus process $\to$ \textbf{6-month cooling} before you can sit GEP again. Offices they hire campus for: \textbf{Mumbai (Airoli / CJB)} and \textbf{Hyderabad (Madhapur)}. Not Noida.

\section{4. Project questions (this is where they actually decide)}
\subsection{LogiFlow --- they will start here}
\Q{Tell me about LogiFlow.}
\Typ{``It is an AI logistics app we made for a hackathon, we used Python and it got Top 106.''}
\Say{We ranked logistics corridors on time, cost, and delay risk so an operator can pick a route under constraints. Google Solution Challenge Global Top 106. I was technical lead on the railway pipeline, the comparator pipeline, and the hybrid pipeline. Data: 15{,}000+ train-day records, 580+ corridors. FastAPI + Next.js UI. Redis for hot reads around 100--400\,ms. pytest 100/100 on pricing/routing rules. Ranking is the product; ML delay signal is one input, not the whole app.}

\Q{What was \emph{your} code vs the team's?}
\Typ{``I did a bit of everything / I was the main lead of the whole product.''}
\Say{Co-lead. I owned railway path: geometry/reads, ranking into the API, tests. Comparator/hybrid: compare modes so the user sees trade-offs, not one dump. Frontend was shared; I do not claim I wrote every React file.}

\Q{Why Redis? What if Redis is down?}
\Typ{``Redis is fast. We used it for speed.''}
\Say{Corridor metadata is read far more than written. Cache-aside: miss $\to$ SQL $\to$ SET with TTL. If Redis is down, FastAPI still answers from SQL, slower. Redis is not the source of truth. Stale reads inside TTL are OK for maps; not OK for an irreversible commit.}

\Q{How do you rank time vs cost vs delay?}
\Typ{``We used AI / a neural network.''}
\Say{Multi-objective. Pareto-style: do not hide a slow cheap route. Delay risk came from a gradient-boosting model on 15k+ train-days --- I can say MAE around 22 minutes if asked, and that it is sklearn-style boosting, not PyTorch. If I am unsure of a number I will not invent it.}

\Q{This is procurement. How is routing relevant?}
\Typ{``GEP and LogiFlow are the same product.''}
\Say{GEP is buy + move + plan. LogiFlow is the \textbf{move} layer: given constraints, which path. I have not built a purchase-order system. I have built decisioning on operational logistics data.}

\subsection{Community Hero}
\Q{What is Community Hero?}
\Typ{``A social app with Gemini.''}
\Say{Civic issue reporting PWA. Vibe2Ship Global Top 20. I solo-owned React/Express, Firestore, Cloud Run. Citizen drops a geo-tagged report; Gemini categorizes; admin reviews. Guard checks: empty/noisy input never hits the model blindly; bad model output does not break routing --- fallback to manual review.}

\Q{What if Gemini is wrong or down?}
\Typ{``Then the app fails / we retry.''}
\Say{Treat the model as untrusted. Validate input. If output is not a known category, store raw text and flag for admin. Prefer a missed auto-tag over a wrong auto-route.}

\subsection{Career Automation Stack (on the resume --- they can open GitHub)}
\Q{What is this stack? The GitHub only has YAML.}
\Typ{``I built 14k questions and emailed 10k recruiters with AI.''}
\Say{Personal scheduled pipelines, not a GEP product. Internship Scout: ingest $\to$ dedup $\to$ HTML digest over 966 companies, GitHub Actions. Hiring Scout: public-source contacts, rate-limited SMTP --- I do not spam. OA Forge: 14k+ evidence-tracked questions, 5.5k+ tests, hybrid C++ judge with g++ and subprocess limits. The public repo is \textbf{workflow YAML only} so personal data stays private. I can walk a job graph: checkout, Python, secrets, cache, commit state. Application code is private; I still own the design.}

\Q{Is scraping / cold email legal? Why should GEP trust you?}
\Typ{``Everyone does it / it is just automation.''}
\Say{Personal learning scope: public pages, rate limits, no purchased lists. At GEP I follow company policy and legal. The engineering I would transfer is scheduling, retries, idempotent state, and tests --- the same discipline as a procurement job that must not double-send.}

\Q{How does the C++ judge work?}
\Say{User code $\to$ compile with g++ (C++17, -O2) $\to$ run under time/memory limits in a subprocess $\to$ compare stdout to expected. If compile fails, verdict is compile error, not WA. I do not claim a kernel sandbox like Codeforces.}

\subsection{Cursor / MCP (only if they bite)}
\Q{What is MCP? Did AI write your resume projects?}
\Typ{``I use all the newest AI tools, I have 20 plugins.''}
\Say{Cursor is my IDE. MCP lets the agent call tools I already use --- I will not list product names unless asked. I write the spec, I review every diff, I merge. Same as a junior who types fast: I still own the code. 28k+ Python LOC and 446 commits in that workflow is volume, not a claim that the agent is the author of record.}

\subsection{Nexus}
\Q{Chairperson --- what do you actually own?}
\Typ{``I manage 500--600 students daily.''}
\Say{Departmental technical cell of CSE and AI. I own planning and delivery of workshops, contests, mentoring --- 500-plus students across events, not a daily headcount. Journey if asked: Member $\to$ Social Media Coordinator $\to$ Academic Mentor $\to$ Technical Representative $\to$ Chairperson, about two years. GEP said they like ownership. This is the people version of the same habit.}

\subsection{IFFCO (on resume --- they may open it even if intro skipped it)}
\Q{What did you do at IFFCO?}
\Typ{``Auto authentication and professional deadlines.''}
\Say{Jun--Jul 2025, technical team. Node.js, Express, MySQL. 10+ REST APIs with validation and error handling. Tools that automated 50+ daily plant workflows. Git, Docker, CI/CD. First time production: logs and stack traces when a flow broke, then a fix, then tell the mentor.}

\section{5. OOP --- 10 minutes (almost every GEP R1)}
\Q{Four pillars, with an example from your code.}
\Typ{Definitions copied from GFG, no example.}
\Say{\textbf{Encapsulation:} hide MySQL password/connection inside a module; handlers only see \texttt{createReport()}. \textbf{Abstraction:} FastAPI route is ``get ranked corridors,'' not the SQL+Redis steps. \textbf{Inheritance:} a base \texttt{Pipeline} with \texttt{run()} / rail vs hybrid override. \textbf{Polymorphism:} comparator and railway both expose rank, different internals. Access: public API, private helpers.}

\Q{Abstract class vs interface.}
\Typ{``Same thing.''}
\Say{Abstract class: shared code + some unimplemented methods. Interface/protocol: only the contract. Python: ABC vs Protocol. I use a shared base when pipelines share validation; a protocol when only the rank() shape is shared.}

\Q{Design a bank / cart with OOP.}
\Typ{Draw random boxes.}
\Say{Classes: Customer, Account, Transaction. Account has withdraw() that checks balance (encapsulation). Savings/Current inherit Account. Transfer is a Transaction. For ``last item in cart'': inventory row with version/quantity; UPDATE \ldots WHERE qty $>$ 0; if 0 rows updated, tell the user. That is optimistic concurrency, not hope.}

\section{6. OS + DBMS + SQL (they fail people here)}
\Q{Process vs thread. Deadlock.}
\Typ{One-line mix-up.}
\Say{Process = program in execution, own memory. Thread = lighter, shares process memory. Deadlock needs all four: mutual exclusion, hold and wait, no preemption, circular wait. Mutex = one owner. Semaphore = counter.}

\Q{ACID. Index. JOIN.}
\Typ{Expand the acronym only.}
\Say{Atomicity all-or-nothing; Consistency rules; Isolation concurrent txs; Durability after commit. Index: faster lookup, slower writes. INNER JOIN: matches both sides. LEFT JOIN: all left + NULL if no match. GROUP BY then HAVING for aggregates. WHERE cannot filter COUNT(*).}

\Q{Write SQL: departments with more than 5 employees.}
\Say{\texttt{SELECT dept, COUNT(*) FROM emp GROUP BY dept HAVING COUNT(*) > 5}}

\Q{SQL vs NoSQL --- you have MySQL and Firestore.}
\Say{LogiFlow/IFFCO: joins, tariffs, validation --- SQL. Community Hero: documents per report, geo, fast iterate --- Firestore. Not ``SQL is old.''}

\Q{MySQL vs MongoDB (TalentBattle GEP). You listed Firestore, not Mongo.}
\Say{I have not used MongoDB in a shipped project. Closest: Firestore documents. SQL when I need joins and constraints. Document store when the record is a blob with nested fields and the schema is still moving.}

\Q{Python vs Java.}
\Say{Python: fast to ship APIs, data, tests (LogiFlow, scouts). Java: typed, JVM, common in enterprise. I have Java in coursework. GEP FTE often uses C\# / Angular --- I have not faked those on the resume; I can learn the product stack.}

\Q{Stack vs queue vs list vs tree.}
\Say{Stack LIFO (undo, DFS, call stack). Queue FIFO (BFS, job queue). List: O(1) head insert, O($n$) index. Tree: hierarchy; BST left $<$ root $<$ right. Graph: corridors in LogiFlow --- nodes + edges, BFS for unweighted hops.}

\section{7. DSA they actually give (Easy--Medium)}
Coding loop every time: constraints, example, brute, better, code, dry-run, O(time)/O(space).

\Q{Reverse an array / string.}
\Say{Two pointers $i,j$; swap until they meet. O($n$) time, O(1) extra.}

\Q{First non-repeating character / remove all repeating characters (GEP intern R1, quoted).}
\Say{Count in a map (or array[256]). First unique: second pass, first count==1. Remove repeats: keep chars with count==1, or use two pointers + map. Say an example: \texttt{swiss} $\to$ first unique \texttt{w}.}

\Q{Two Sum / frequency map.}
\Say{Hash map. Two Sum: store complement. O($n$).}

\Q{Middle of a linked list.}
\Say{Slow/fast pointers. Fast moves 2, slow moves 1. When fast hits end, slow is mid. O($n$) time, O(1) extra.}

\Q{Alt+Tab OA (quoted on GFG Associate SD): $n=4$, list $\{1,2,3,4\}$, $k=3$ $\to$ $\{3,1,2,4\}$. Catch $k \ge n$.}
\Say{The $k$-th window comes to front. Index $i=(k-1)\bmod n$; rotate that element to index 0, shift the prefix right. Dry-run $k=3$ and $k=7$ before coding.}

\Q{Seconds between two clock times (GEP intern OA, easy).}
\Say{Parse HH:MM:SS to seconds since midnight; subtract; if negative, add 24h. Watch midnight wrap.}

\Q{Kadane --- max subarray. Dry-run \texttt{[-2,1,-3,4,-1,2,1,-5,4]}.}
\Say{Running sum; if sum $<$ 0 reset to 0; track max. Answer 6, subarray $[4,-1,2,1]$.}

\Q{Binary search.}
\Say{\texttt{lo=0, hi=n-1; while lo<=hi: mid=(lo+hi)//2}. Sorted only. O($\log n$).}

\Q{Fibonacci in a range (GFG GEP R2).}
\Say{Iterate $a,b=0,1$; print while in range. Two methods: loop and formula if they want. Watch overflow.}

\Q{BFS / trees (Sarthak's GEP intern R2).}
\Say{Favourite DS if asked: \textbf{hash map} for APIs or \textbf{graph/queue} for corridors --- pick one and stay. BFS: queue, level-order, shortest unweighted path. Level-order: push root; pop; push left then right.}

\Q{Primes 1 to 100.}
\Say{Trial division or sieve. Sieve: mark multiples. Do not say ``I used AI to generate sieve'' --- write it.}

\section{8. Puzzles (GFG GEP R2)}
\Q{12 balls, 1 different weight, 3 weighings.}
\Say{Ternary split 4--4. Balance tells which group. Repeat. Classic.}

\Q{3L and 5L jugs, make 4L.}
\Say{Fill 5, pour into 3 (2 left in 5), empty 3, pour 2 into 3, fill 5, pour into 3 (1 space) $\to$ 4 in the 5L.}

If stuck: state assumptions, try a small case, do not freeze.

\section{9. AI they may ask (PPT: wide idea, not research)}
\Q{AI vs ML vs DL. Supervised vs unsupervised.}
\Say{AI broad; ML learns from data; DL neural nets. Supervised = labels (classify/regress). Unsupervised = clusters. Linear regression predicts a number, not a class.}

\Q{How would you put an LLM on procurement data?}
\Typ{``LangGraph agents, RAG, fine-tune GPT.''}
\Say{I have not used LangGraph. Pattern I have used: ingest $\to$ validate $\to$ model (Gemini categorize) $\to$ guard $\to$ human/admin if unsure. Same shape on invoices/POs: extract fields, never auto-post a bad total. I would learn their workflow tool on the job.}

\Q{Hadoop on OA (seen on some AI-intern / SE papers).}
\Say{Distributed processing of large files on a cluster. Map then reduce. I have \textbf{not} run a Hadoop job. I will not pretend. SQL/Python pipelines are what I have shipped.}

\Q{Purpose of OOP / DS vs algorithms (short intern R2).}
\Say{OOP: bundle data+behavior so a CorridorRanker can change internals without breaking callers. DS = how we store (hash, tree). Algorithms = how we process (BFS, Kadane). You pick a DS to make the algorithm cheap.}

\Q{How do you know a website's language without reading the code?}
\Say{Look at \texttt{lang} on HTML, \texttt{Content-Language} header, or script URLs. If they want CS: detect from tokens/stopwords. I say I am guessing and ask which layer they mean --- HTTP, HTML, or NLP.}

\section{10. HR --- typical questions and answers}
HR is short. People fail by lying on location, badmouthing college, or having no reason for GEP except stipend.

\Q{What do you know about GEP? (intern HR, quoted)}
\Typ{Recite Wikipedia or name random modules.}
\Say{Software + consulting + managed services for procurement and supply chain. Products I will name at a high level: SMART (source-to-pay), NEXXE (supply chain), Quantum Intelligence as the AI layer. I will not fake screens. Tech interviews almost never quiz S2P cycles --- that is the consulting track.}

\Q{Why GEP?}
\Typ{``Great brand, good package, work-life balance.''}
\Say{AI-powered procurement and supply chain --- software plus how enterprises actually buy and move. LogiFlow is routing decisions on operational data. I want that class of problem at company scale, with ownership.}

\Q{Where do you see yourself in 5 years? (intern R2/HR)}
\Say{Engineer who can own a service on their procurement/supply-chain stack --- APIs, data, tests --- and who can add AI behind guards. Not ``CTO.''}

\Q{Hyderabad or Mumbai?}
\Typ{Pick one then flip later / ``wherever, I don't care.''}
\Say{I am open to both. No family constraint. I will go where the team is. (Do not invent a favourite.)}

\Q{Why software if your degree is AI?}
\Typ{``AI is just a name / I want to switch.''}
\Say{I ship systems --- APIs, data, tests. AI is a component (ranking, Gemini) on top of that, with guards. GEP needs both.}

\Q{Strength / weakness.}
\Typ{``Perfectionist.''}
\Say{Strength: I finish --- design to deploy, I own it (Nexus + Community Hero). Weakness: I can over-build a personal tool. I time-box and ship a thin version first, then iterate --- that is how I keep scout pipelines from never releasing.}

\Q{Family background / where are you from?}
\Say{Short, factual. Prayagraj roots, college in Surat. No drama.}

\Q{Do you have offers? Will you join?}
\Say{If true: this is the process I am in; if I get GEP I intend to join. Do not play companies against them.}

\Q{Questions for us?}
\Typ{``No'' or ``What is the stipend?''}
\Say{How do interns contribute to AI / agent workflows on procurement data in the first month? What does a good intern look like at 8 weeks?}

\section{11. Weak vs strong --- same question}
\noindent\begin{tabularx}{\textwidth}{@{}>{\bfseries}p{3.2cm} >{\color{muted}}X >{\color{ok}}X@{}}
Topic & Weak (common) & Strong (you) \\
Intro & 3 minutes, IFFCO story, location, 7 MCP names & Locked 50s intro, then stop \\
LogiFlow & ``AI app, I did everything'' & Rank time/cost/delay; three pipelines; tests \\
Gemini & ``We used AI for intelligence'' & Guard + fallback; model is untrusted \\
DSA & Jump to code, silent & Clarify, brute, optimize, dry-run \\
Don't know & Guess LangGraph / Azure & ``I have not used X; closest is Y'' \\
HR location & Only Mumbai / only Hyd & Both, no constraint \\
\end{tabularx}

\section{12. Day-of rules}
\begin{itemize}
  \item Print the submitted resume. Do not explain a project that is not on it (AirHelp: only if they ask GitHub).
  \item Every noun you say is a follow-up: Redis, pytest, Firestore, Gemini, MCP, Chairperson.
  \item If coding: talk. Partial working brute force $>$ silent optimal.
  \item R1: no stipend, no location speech. R2/HR: location both.
  \item End each round: ``Thank you. I enjoyed this.'' Then one question if they offer time.
\end{itemize}

\section{13. Sources (so you know this is not invented)}
\begin{itemize}
  \item GEP Campus Connect India: 2-day close, Hyd + Mumbai, intern for 2028 passout --- \texttt{gep.com/careers/join-us/campus-connect/india}
  \item GFG: GEP SDE --- OA (aptitude, maths, OS/DBMS, 1 coding); R1 OOP+resume DS+project+array+puzzle; R2 project ML + Fibonacci in range; HR family + location
  \item GFG: GEP SDE-1 on-campus --- R2 OOP bank design + last-item cart + 12 balls + jugs
  \item LinkedIn: Shivam Kumar (GEP intern $\to$ SDE 2025) --- 2 tech + HR; 1 Easy--Medium DSA/round; questions from intro/resume
  \item LinkedIn: Sarthak Chauhan (SVNIT, GEP intern) --- OOP R1; trees+BFS R2; HR relocate
  \item AmbitionBox: ``Explain key projects you implemented''; GEP SWE stack often C\# / Angular / SQL for FTE --- you do not fake that; you say open to learn
  \item GFG intern on-campus: first unique char; difficulties in project; short R2; everyone who reached HR got intern
  \item TalentBattle ASE 2023: OOP, access specifiers, stack/queue/list/tree, mid of LL, BS, primes 1--100, MySQL vs Mongo
  \item UnsaidTalks: Maanya Jain AI intern 2024 (Preliminary + Functional + HR; Hadoop on OA); Sanjay Gupta (intern Mumbai, FTE Hyd or Mumbai)
  \item Official cooling period 6 months on fail. Campus offices: Mumbai CJB and Hyderabad only.
  \item Your PPT: two interviews; ownership; wide AI; supply chain. Your submitted PDF is the only source of truth in the room.
\end{itemize}

\vspace{2mm}
\noindent{\footnotesize\color{hot}\textbf{Golden rule:} If it is not on the resume and you cannot dry-run it, do not say it. Clarify $\to$ approach $\to$ code $\to$ dry-run $\to$ complexity. Never silent.}
\end{document}
